% Part: turing-machines
% Chapter: machines-computations
% Section: disciplined-machines

\documentclass[../../../include/open-logic-section]{subfiles}

\begin{document}

\olfileid{tur}{mac}{dis}
\olsection{الآلات المنضبطة}

\begin{explain}
نظرنا في القسم \olref[hal]{sec} في آلات تورنغ ذات حالة توقف واحدة
معينة~$h$؛ ومضمون لها، إن توقفت أصلًا، أن تتوقف في الحالة~$h$. وبهذا
المعنى تكون الآلات ذات حالة التوقف الواحدة «أشد انضباطًا» مما نسمح به
لآلات تورنغ بوجه عام. وثمة قيود أخرى يمكن أن نفرضها على سلوك آلات
تورنغ. فمثلًا لم نحظر عليها محو علامة نهاية الشريط من الخانة~$0$، ولا
محاولة التحرك يسارًا من الخانة~$0$. (ينص تعريفنا على أن الرأس يبقى
ببساطة في الخانة~$0$ في هذه الحال؛ أما تعريفات أخرى فتجعل الآلة
تتوقف.) ومن المرغوب فيه أحيانًا أيضًا أن نستطيع افتراض أن آلة تورنغ،
إن توقفت أصلًا، تتوقف في الخانة~$1$.
\end{explain}

\begin{defn}\ollabel{defn:disciplined}
تكون آلة تورنغ~$M$ \emph{منضبطة} إذا وفقط إذا
\begin{enumerate}
    \item كانت لها حالة توقف واحدة معينة~$h$،
    \item توقفت، إن توقفت أصلًا، وهي تمسح الخانة~$1$،
    \item لم تمح قط الرمز $\TMendtape$ من الخانة~$0$، و
    \item لم تحاول قط التحرك يسارًا من الخانة~$0$.
\end{enumerate}
\end{defn}

\begin{explain}
سبق أن ناقشنا إمكان تحويل كل آلة تورنغ إلى آلة لها السلوك نفسه وحالة
توقف معينة. ويتم ذلك ببساطة بإضافة حالة جديدة~$h$، وإضافة تعليمة
$\delta(q,\sigma)=\tuple{h,\sigma,\TMstay}$ لكل زوج
$\tuple{q,\sigma}$ تكون عنده الدالة الأصلية~$\delta$ غير معرّفة.
وصحيح، وإن كان إثباته مضجرًا، أنه يمكن تحويل كل آلة تورنغ~$M$ إلى آلة
تورنغ منضبطة~$M'$ تتوقف عند المدخلات نفسها وتنتج الخرج نفسه. فمثلًا،
إذا توقفت آلة تورنغ ولم تكن في الخانة~$1$، أمكننا إضافة تعليمات تجعل
الرأس يتحرك يسارًا حتى يجد علامة نهاية الشريط، ثم يتحرك خانة واحدة
نحو اليمين، ثم يتوقف. وسنترك لك التفكير في كيفية معالجة الشروط
الأخرى.
\end{explain}

\begin{ex}
    في \olref{fig:adder-disc} نحول آلة الجمع الواردة في
    \olref[una]{ex:adder} إلى آلة منضبطة.
    \begin{figure}\[
\begin{tikzpicture}[->,>=stealth',shorten >=1pt,auto,node distance=2.8cm,
                    semithick]
  \tikzstyle{every state}=[fill=none,draw=black,text=black]

  \node[initial,state]         (A)              {$q_0$};
  \node[state]         (B) [right of=A] {$q_1$};
  \node[state]         (C) [below right of=B] {$q_2$};
  \node[state]         (D) [below left of=C] {$q_3$};
  \node[state]         (H) [left of=D] {$h$};

  \path (A) edge node {\TMtrans{\TMblank}{\TMstroke}{\TMstay}} (B)
           edge [loop below] node {\TMtrans{\TMstroke}{\TMstroke}{\TMright}} (A)
        (B) edge [loop below] node {\TMtrans{\TMstroke}{\TMstroke}{\TMright}} (B)
            edge node {\TMtrans{\TMblank}{\TMblank}{\TMleft}} (C)
        (C) edge node {\TMtrans{\TMstroke}{\TMblank}{\TMleft}} (D)
        (D) edge [loop above] node {\TMtrans{\TMstroke}{\TMstroke}{\TMleft}} (D)
        edge node {\TMtrans{\TMendtape}{\TMendtape}{\TMright}} (H);
\end{tikzpicture}
\]\caption{آلة جمع منضبطة}
\ollabel{fig:adder-disc}
\end{figure}
\end{ex}

\begin{prop}\ollabel{prop:disciplined} لكل آلة تورنغ~$M$ آلة تورنغ
منضبطة~$M'$ تتوقف بخرج~$O$ إذا توقفت $M$ بخرج~$O$، ولا تتوقف إذا لم
تتوقف~$M$. وعلى وجه الخصوص، كل دالة $f\colon\Nat^n\to\Nat$ تحسبها
آلة تورنغ تحسبها أيضًا آلة تورنغ منضبطة.
\end{prop}

\begin{prob}\label{tur:mac:dis:prob:disc-succ}
أعط آلة منضبطة تحسب $f(x)=x+1$.
\end{prob}


\begin{prob}\label{tur:mac:dis:prob:copier}
جد آلة منضبطة تنتج، إذا بدأت عند الدخل $\TMstroke^n$، الخرج
$\TMstroke^n\concat\TMblank\concat\TMstroke^n$.
\end{prob}

\end{document}
